\documentclass[journal]{IEEEtran}
\usepackage[utf8]{inputenc}
\usepackage[T1]{fontenc}
\usepackage[english]{babel}
\usepackage{cite,graphicx,array,amsmath,booktabs,needspace,placeins}
\usepackage[hidelinks]{hyperref}
\usepackage{listings,xcolor}
\lstdefinestyle{arduino}{
    language=C++,
    basicstyle=\ttfamily\scriptsize,
    keywordstyle=\color{blue!70!black},
    commentstyle=\color{green!40!black},
    stringstyle=\color{red!60!black},
    numbers=left,
    numberstyle=\tiny\color{gray},
    stepnumber=1,
    numbersep=5pt,
    breaklines=true,
    breakatwhitespace=false,
    frame=single,
    rulecolor=\color{gray!50},
    backgroundcolor=\color{gray!5},
    showstringspaces=false,
    tabsize=2,
    columns=fullflexible,
    keepspaces=true,
    captionpos=b,
    xleftmargin=0pt,
    xrightmargin=0pt
}


\newcommand{\labpath}{}
\IfFileExists{4/introduction.tex}{\renewcommand{\labpath}{4/}}{}
\graphicspath{{\labpath img/}}
\begin{document}
\raggedbottom
\title{LAB 4: LoRa P2P and Arduino MEGA Integration}
\author{José~Francisco~Paca~Sotero, Daniel~Ignacio~Casquino~Paz,\\
Nicolas~Stigler~Yañez, Sofia~Valentina~Ku~Paredes\\
Universidad de Ingeniería y Tecnología}
\maketitle
\section{Introduction}
\input{\labpath introduction.tex}
\section{Methodology}
\input{\labpath methodology.tex}
\section{Development and Discussion}
\input{\labpath development_discussion.tex}
\FloatBarrier
\Needspace{12\baselineskip}
\section{Conclusions and Recommendations}
\input{\labpath conclusion.tex}
\newpage
\bibliographystyle{IEEEtran}
\bibliography{\labpath references}
\end{document}
